\documentclass[11pt]{article}
\usepackage[margin=1in]{geometry}
\usepackage{enumitem}
% =============================================================================
% Preambles/header.tex — shared LaTeX/Beamer preamble for this project
%
% Use from a Beamer lecture by prepending:
%     \input{header}
% and compiling with TEXINPUTS=../Preambles:$TEXINPUTS (the /compile-latex
% skill does this automatically).
%
% PALETTE CONTRACT. Color names defined here MUST match the names in
% Quarto/theme-template.scss so Beamer and Quarto renderings use the same
% palette. The scripts/check-palette-sync.sh script enforces this.
%
% To customize: edit the HEX values below AND the matching $variable in
% Quarto/theme-template.scss, then run scripts/check-palette-sync.sh.
% =============================================================================

% -----------------------------------------------------------------------------
% Engine detection + encoding
%
% Our /compile-latex skill uses XeLaTeX, where inputenc is unnecessary and
% can produce encoding warnings. Only load inputenc under pdfTeX.
% -----------------------------------------------------------------------------
\usepackage{iftex}
\ifPDFTeX
  \usepackage[utf8]{inputenc}
\fi

\usepackage{amsmath, amssymb, amsthm}
\usepackage{booktabs}
\usepackage{graphicx}

% -----------------------------------------------------------------------------
% PALETTE — mirrors Quarto/theme-template.scss variable names.
% Load xcolor and define colors BEFORE anything that references them
% (notably hyperref, hypersetup, and the Beamer color assignments).
% -----------------------------------------------------------------------------
\usepackage{xcolor}

\definecolor{primary-blue}{HTML}{012169}      % headings, accents
\definecolor{primary-gold}{HTML}{B9975B}      % emphasis, borders
\definecolor{highlight-yellow}{HTML}{F2A900}  % markers, alerts
\definecolor{light-bg}{HTML}{E8EDF5}          % title / section backgrounds
\definecolor{jet}{HTML}{1A1A1A}               % body text

% Semantic colors (used in diagrams and annotations)
\definecolor{positive}{HTML}{15803D}          % good / observed / identified
\definecolor{negative}{HTML}{B91C1C}          % bad / problematic / bias
\definecolor{neutral}{HTML}{525252}           % reference / context

% Highlight accents (match .hi-* classes in the SCSS)
\definecolor{hi-slate}{HTML}{314F4F}
\definecolor{hi-green}{HTML}{2E8B57}
\definecolor{hi-red}{HTML}{C41E3A}

% Semantic aliases so skill templates and snippets can write
% \textcolor{accent}{...} without hard-coding "primary-blue".
\colorlet{accent}{primary-blue}
\colorlet{accent2}{primary-gold}

% -----------------------------------------------------------------------------
% Hyperref — loaded AFTER xcolor + palette so link colors resolve.
% -----------------------------------------------------------------------------
\usepackage{hyperref}
\hypersetup{colorlinks=true, linkcolor=primary-blue, urlcolor=primary-blue,
            citecolor=primary-blue, pdfborder={0 0 0}}

% -----------------------------------------------------------------------------
% Course code — set once per deck (after \input{header}, before \begin{document})
% via \coursecode{CS401}. Shown in the footer on every slide so a deck's course
% is identifiable without opening the title page. Empty by default.
% -----------------------------------------------------------------------------
\makeatletter
\newcommand{\coursecode}[1]{\gdef\@coursecodetext{#1}}
\gdef\@coursecodetext{}
\makeatother

% -----------------------------------------------------------------------------
% Beamer theme assignments (only apply under Beamer).
%
% We detect Beamer via \@ifclassloaded{beamer}, which is the documented,
% non-internal way. This must be wrapped in \makeatletter/\makeatother
% because \@ifclassloaded contains an @-catcode character.
% -----------------------------------------------------------------------------
\makeatletter
\@ifclassloaded{beamer}{%
  \usefonttheme{professionalfonts}
  \setbeamercolor{structure}{fg=primary-blue}
  \setbeamercolor{frametitle}{fg=primary-blue}
  \setbeamercolor{title}{fg=primary-blue}
  \setbeamercolor{subtitle}{fg=primary-gold}
  \setbeamercolor{author}{fg=primary-blue}
  \setbeamercolor{institute}{fg=neutral}
  \setbeamercolor{date}{fg=primary-gold}
  \setbeamercolor{itemize item}{fg=primary-blue}
  \setbeamercolor{itemize subitem}{fg=primary-gold}
  \setbeamercolor{alerted text}{fg=primary-gold}
  \setbeamercolor{block title}{fg=white, bg=primary-blue}
  \setbeamercolor{block body}{bg=light-bg}
  % Semantic block variants: block=definition/concept (blue), exampleblock=
  % worked example (green/positive), alertblock=key takeaway or caution (gold).
  % Keep to <=2 colored boxes per slide across ALL three variants (INV-7).
  \setbeamercolor{block title example}{fg=white, bg=positive}
  \setbeamercolor{block body example}{bg=light-bg}
  \setbeamercolor{block title alerted}{fg=white, bg=primary-gold}
  \setbeamercolor{block body alerted}{bg=light-bg}

  % Minimal footer: course code (left) + page X / N (right), no navigation clutter
  \setbeamertemplate{navigation symbols}{}
  \setbeamertemplate{footline}{%
    \color{neutral}\footnotesize\hspace{0.7em}\@coursecodetext\hfill
    \insertframenumber/\inserttotalframenumber\hspace{0.7em}\vspace{0.3em}}
}{}
\makeatother

% -----------------------------------------------------------------------------
% TikZ setup — libraries the snippet gallery uses
% -----------------------------------------------------------------------------
\usepackage{tikz}
\usetikzlibrary{arrows.meta, positioning, calc, decorations.pathreplacing,
                fit, shapes.geometric, backgrounds}

% Shared TikZ styles — snippets and hand-written diagrams can pick these up
% via `\begin{tikzpicture}[every node/.style={...}]` or by naming specific
% styles (e.g., `\node[dag-node] (X) at (0,0) {X};`).
\tikzset{
  % Node styles for causal DAGs and similar
  dag-node/.style = {
    draw=primary-blue, fill=light-bg, rounded corners=2pt,
    minimum width=1.2cm, minimum height=0.9cm, align=center, font=\small
  },
  decision-node/.style = {
    draw=primary-gold, fill=white, diamond, aspect=1.8,
    minimum width=1.8cm, minimum height=1.2cm, align=center, font=\small,
    inner sep=1pt
  },
  % Edge styles — observed vs counterfactual
  observed-edge/.style = {-{Latex[length=2.5mm]}, thick, draw=jet},
  counterfactual-edge/.style = {-{Latex[length=2.5mm]}, thick, draw=neutral, dashed},
  confound-edge/.style = {-{Latex[length=2.5mm]}, thick, draw=negative, dashed},
  % Data-point styles for plots
  observed-dot/.style = {fill=primary-blue, draw=primary-blue, circle, inner sep=1.2pt},
  counterfactual-dot/.style = {fill=white, draw=neutral, circle, inner sep=1.2pt}
}

% -----------------------------------------------------------------------------
% Convenience macros
% -----------------------------------------------------------------------------
\newcommand{\muted}[1]{\textcolor{neutral}{#1}}
\newcommand{\key}[1]{\textcolor{primary-gold}{\textbf{#1}}}
\newcommand{\good}[1]{\textcolor{positive}{#1}}
\newcommand{\bad}[1]{\textcolor{negative}{#1}}

% Standout frame macro: full-bleed dark background for section transitions.
% Beamer-only (uses \begin{frame}/\setbeamercolor/\usebeamerfont) -- do not
% call from an article-class file (e.g. Notes/<CODE>/*-notes.tex). \muted,
% \key, \good, \bad, and \coursecode above are plain xcolor/gdef and are
% safe to use from either class; verified when Notes/article-class support
% was added.
% Expand: \transitionslide{Your transition title}
\newcommand{\transitionslide}[1]{%
  \begingroup
  \setbeamercolor{background canvas}{bg=primary-blue}
  \begin{frame}[plain]
    \centering
    \vfill
    {\usebeamerfont{title}\color{white}\Large #1\par}
    \vfill
  \end{frame}
  \endgroup
}

% Section divider frame (Beamer-only), an alternative to \transitionslide with a
% darker two-line style: near-black (jet) background, a small white label line
% above a larger gold title, and a thin gold rule along the bottom edge.
% Expand: \sectiondivider{Section 2}{Pointers}
\newcommand{\sectiondivider}[2]{%
  \begingroup
  \setbeamercolor{background canvas}{bg=jet}
  \begin{frame}[plain]
    \vspace*{3.0cm}
    {\color{white}\ttfamily\large #1\par}
    \vspace{0.5cm}
    {\color{primary-gold}\LARGE #2\par}
    \begin{tikzpicture}[remember picture, overlay]
      \draw[primary-gold, line width=2.5pt]
        ($(current page.south west)+(0,0.1)$) -- ($(current page.south east)+(0,0.1)$);
    \end{tikzpicture}
  \end{frame}
  \endgroup
}

% -----------------------------------------------------------------------------
% Channel watermark — "Concepts That Click" (YouTube).
%
% Article-class files (CompetitiveExam Q/A sets, Notes, Assignments) get a
% light diagonal draft watermark on every page plus a centered footer naming
% the channel. Beamer decks get a subtle TikZ background watermark instead
% (the footline already carries the course code). Centralized here so every
% MCQ PDF is branded without per-file edits.
% -----------------------------------------------------------------------------
\makeatletter
\@ifclassloaded{article}{%
  \usepackage{draftwatermark}
  \SetWatermarkText{Concepts That Click}
  \SetWatermarkScale{1}
  \SetWatermarkFontSize{2cm}
  \SetWatermarkLightness{0.86}
  \SetWatermarkAngle{45}
  \usepackage{fancyhdr}
  \pagestyle{fancy}
  \fancyhf{}
  \fancyfoot[C]{\footnotesize\textcolor{neutral}{Concepts That Click~|~YouTube: \href{https://www.youtube.com/@concepts.thatclick}{@concepts.thatclick}}}
  \renewcommand{\headrulewidth}{0pt}
  \renewcommand{\footrulewidth}{0pt}
  \fancypagestyle{plain}{%
    \fancyhf{}%
    \fancyfoot[C]{\footnotesize\textcolor{neutral}{Concepts That Click~|~YouTube: \href{https://www.youtube.com/@concepts.thatclick}{@concepts.thatclick}}}%
    \renewcommand{\headrulewidth}{0pt}%
    \renewcommand{\footrulewidth}{0pt}%
  }
}{}
\@ifclassloaded{beamer}{%
  \setbeamertemplate{background}{%
    \begin{tikzpicture}[remember picture, overlay]
      \node[rotate=45, opacity=0.055, align=center]
        at (current page.center)
        {\Huge\textcolor{neutral}{Concepts That Click}};
    \end{tikzpicture}%
  }
}{}
\makeatother 
\coursecode{JKSSB}
\renewcommand{\thesection}{17.\arabic{section}}
\newtheorem{example}{Example}[section]
\newenvironment{definitionbox}[1]{%
  \par\noindent\textbf{Definition (#1).}\ \itshape
}{\par\vspace{0.3em}}
\title{Hardware, Software and Firmware --- Detailed Lecture Notes}
\author{JKSSB Computer Awareness}
\date{\today}

\begin{document}
\maketitle

\section{Hardware}
Hardware is the set of physical, tangible parts of a computer system --- the
parts you can see and touch. A computer's hardware includes the central
processing unit (CPU), the motherboard, random access memory (RAM), the hard
disk drive or solid-state drive, the keyboard, the mouse, the monitor, the
printer, and every other physical component. Hardware carries out the actual
work of the machine: it executes instructions, stores data, and moves signals.

Hardware alone does nothing useful. A processor with no instructions to run
is only silicon and metal. The instructions come from software, and the two
together make a working computer. For exam purposes, the test for hardware is
simple: if it is a physical object that you can touch, it is hardware.

\begin{definitionbox}{Hardware}
The physical equipment of a computer system --- the tangible components such
as the CPU, memory, storage, input devices, output devices, and communication
hardware.
\end{definitionbox}

Hardware is usually grouped by function. Input devices, such as the keyboard,
mouse, and scanner, bring data into the computer. Output devices, such as the
monitor, printer, and speaker, present results. Processing hardware, mainly
the CPU and the graphics processing unit (GPU), performs computation. Memory
hardware, such as RAM, ROM, and cache, holds data and instructions. Storage
hardware, such as the hard disk drive (HDD) and solid-state drive (SSD),
keeps files without power. Communication hardware, such as the network
interface card and the motherboard, connects the parts together.

Hardware is also described as internal or external. \textbf{Internal
hardware} is fixed inside the computer case --- the motherboard, CPU, RAM, and
internal drives. \textbf{External hardware}, also called a peripheral, is
connected from outside, such as the keyboard, mouse, monitor, printer, and an
external drive. A question may ask whether a named part is internal or
external; the touch test still applies, because both kinds are physical and
both are hardware.

A single concrete machine runs through this course: the \textbf{office
desktop} --- a personal computer with a CPU, memory, a disk, a keyboard, a
mouse, a monitor, and a printer, running an operating system and MS Word.
Every term in this lesson can be pointed at on that one machine, which makes
the abstract layers easier to remember.

\section{Software}
Software is the set of programs and instructions that tell the hardware what
to do. Unlike hardware, software is intangible: you cannot touch it. It is
kept on secondary storage, such as a hard disk or solid-state drive, and it
is loaded into main memory when it runs. Software is what turns a collection
of physical parts into a usable computer.

\begin{definitionbox}{Software}
The collection of programs, procedures, and related data that direct the
operation of a computer. Software is logical and intangible, and it instructs
the hardware.
\end{definitionbox}

Software falls into two broad kinds. \textbf{System software} runs and manages
the computer itself --- the operating system, device drivers, and utility
programs. \textbf{Application software} helps the user perform a task --- a
word processor such as MS Word, a web browser, a spreadsheet, or a game. Every
program a candidate meets in JKSSB questions belongs to one of these two
kinds, but the boundary matters: the operating system is system software, not
application software.

Software is kept on secondary storage as files. When it is launched, the
operating system copies it into main memory, and the CPU executes it from
there. This is why a program must be \emph{loaded} before it can run, and why
data held only in RAM is lost when the power goes off. The permanent copy
lives on disk; the running copy lives in memory.

\section{Hardware vs software}
The two are easy to confuse only when the wording is careless, so hold the
contrasts in a table.

\begin{center}
\begin{tabular}{p{0.18\textwidth}p{0.36\textwidth}p{0.36\textwidth}}
\toprule
\textbf{Point} & \textbf{Hardware} & \textbf{Software} \\
\midrule
Nature & Physical and tangible & Logical and intangible \\
Can be touched? & Yes & No \\
Role & Executes the work & Gives the instructions \\
Stored or fixed & Fixed as physical parts & Stored as files, loaded to run \\
Example & CPU, printer, keyboard & MS Word, a driver, an OS \\
How it fails & Wears out, breaks, burns out & Bugs, corruption, wrong logic \\
\bottomrule
\end{tabular}
\end{center}

Neither can work alone. Hardware and software together form a computer system,
and a question that names one and describes the other is usually testing this
single distinction. A candidate should be able to sort a mixed list --- a CPU,
MS Word, a printer driver, the BIOS, and a user --- into its correct class at a
glance.

\begin{example}
A question lists ``CPU, MS Word, printer driver, BIOS, operator'' and asks
which are software. \textbf{MS Word} and the \textbf{printer driver} are
software. The CPU is hardware; the BIOS is firmware (software stored in
hardware); and the operator is a person, that is, humanware. Sorting is done
by asking, for each item, whether it is physical, a program, embedded code, or
a person.
\end{example}

\section{System software and application software}
\textbf{System software} manages the machine and provides services to other
programs. The operating system is the main example, together with device
drivers and system utilities. \textbf{Application software} is written to let
the user accomplish a task, and it depends on the system software beneath it.

For the exam, the stable rule is: the operating system is system software.
This is the point tested by the real Junior Assistant 2026 concept recorded as
PYQ-06 (device driver = software; OS = system software). A device driver, too,
is system software --- it is a program, not a physical part.

\section{Firmware}
Firmware is software that is embedded in hardware and stored in non-volatile
memory. It is held in read-only memory (ROM), erasable programmable ROM
(EPROM), or flash memory inside the device, so it is available the moment the
device is powered on. Firmware provides the low-level control instructions a
device needs and starts it up.

\begin{definitionbox}{Firmware}
Software permanently stored in a hardware device, usually in non-volatile
memory such as ROM or flash. It is software by nature but is treated as part
of the hardware because it is fixed to a specific device.
\end{definitionbox}

Common examples of firmware include the \textbf{BIOS} (Basic Input Output
System) or \textbf{UEFI} on a personal computer, which initialises the
hardware and begins loading the operating system; the firmware inside a
printer, a router, a keyboard, or a solid-state drive. Firmware is usually
described as sitting between hardware and software. It is a program, so it is
software; it is stored on the physical device, so it behaves like part of the
hardware.

Firmware differs from ordinary software in how tightly it is bound to the
device. Firmware is written for one specific piece of hardware, is stored in
ROM or flash on that device, and changes rarely (it is updated by a process
called flashing). Ordinary software is general-purpose, is stored on secondary
storage, and is installed and updated frequently.

Firmware is the \textbf{first software to run} when a device is switched on. On
a personal computer, the CPU executes the BIOS or UEFI firmware from ROM before
any operating system is available; that firmware checks the hardware and then
hands control to the operating system. This is exactly why firmware must be
non-volatile: it has to be present even before anything can be loaded. A
candidate should not confuse firmware with the operating system --- the
firmware runs first and \emph{starts} the computer, while the operating system
is what the machine is finally running.

\section{Device drivers}
A device driver is a small program that lets the operating system communicate
with a piece of hardware. It acts as a translator: the operating system sends
a general request such as ``print this page'', and the driver converts it into
the exact signals that a specific printer understands. A printer driver, for
example, lets the operating system print on one particular printer model.

\begin{definitionbox}{Device driver}
A program that enables the operating system to communicate with, and control,
a specific hardware device by translating general commands into
device-specific signals.
\end{definitionbox}

The single most important exam point about a driver is that it is
\textbf{software}, even though it is tied to hardware and is often named after
a device. A candidate may be tempted to call a driver hardware because it
belongs to a printer or a keyboard, but that is wrong. The driver is a program,
it is stored like software, and a device will not work properly without the
right driver installed. This is TRAP-10, and it is also the heart of the real
Junior Assistant 2026 concept PYQ-06.

\begin{example}
A question states, ``A device driver is a hardware component that connects the
printer to the computer.'' This is wrong on two counts: the driver is a
\textbf{program}, not a physical component, and it is software, not hardware.
The physical cable and port are hardware; the driver that makes the printer
work is software. The trap is that the driver carries the device's name.
\end{example}

\section{Humanware}
Humanware, also called \textbf{peopleware}, is the human component of a
computer system. It covers everyone involved: the end users who operate the
machine, the programmers who write software, the system administrators who
maintain it, and the operators who run it. A computer system is usually
described as hardware, software, and humanware together.

\begin{definitionbox}{Humanware}
The people who design, operate, maintain, and use a computer system; the human
component that hardware and software ultimately serve.
\end{definitionbox}

Humanware is not a device and not a program. In a classification question, if
the named item is a person or a role --- user, programmer, operator,
administrator --- it is humanware. It is the layer that makes decisions and
uses the results, and without it the hardware and software would have no
purpose.

\section{The layered computing stack}
The parts of a computer system can be arranged as a stack of layers, each one
resting on the layer below. The order, from the bottom up, is:

\begin{center}
\begin{tabular}{p{0.30\textwidth}p{0.58\textwidth}}
\toprule
\textbf{Layer} & \textbf{Role} \\
\midrule
Hardware & The physical components at the base of the system \\
Firmware & Stored on the hardware; starts it and controls it at a low level \\
Device driver & Lets the operating system communicate with the hardware \\
Operating system & Manages all resources and provides services to programs \\
Application software & Lets the user perform a task \\
\bottomrule
\end{tabular}
\end{center}

Reading the stack from the bottom, the physical parts come first. Firmware sits
directly on them and brings the hardware to life. The device driver sits above
the firmware and gives the operating system a clean way to talk to each device.
The operating system manages memory, files, processes, and devices, and offers
services to the programs above it. The application software at the top is what
the user actually opens and works in. A command travels \emph{down} the stack
--- from an application, through the operating system and the driver, to the
hardware --- and results travel back up.

This stack is the single mental picture for the whole lesson: hardware at the
bottom, software above it, and firmware and drivers as the thin bridges that
connect the two.

\section{Why the stack matters}
The layered stack is not only a way to sort terms; it explains the order in
which work is done. When a user clicks Print in a word processor, the command
begins at the application-software layer. The application asks the operating
system to print; the operating system passes the job to the device driver for
that printer; the driver translates the request into the signals the printer
hardware understands; and the printer's own firmware controls the physical
printing. The result then travels back up the stack to the user. Every layer
depends on the one below it, and a fault in any single layer can stop the
whole chain --- which is why ``reinstall the driver'' is such a common fix.

\begin{example}
A printer fails to print although the paper and ink are fine. Because the
system has layers, the fault could lie in the application (a wrong page range),
the operating system (a stuck print queue), the driver (missing or outdated),
or the hardware (a loose cable). Testing from the bottom up checks each layer
in turn. In exam terms, the layer that ``lets the operating system talk to the
printer'' is the \textbf{device driver}.
\end{example}

\section{Exam retrieval and common confusions}
Five distinctions prevent most errors on this topic.
\begin{enumerate}
  \item Hardware is \textbf{physical and tangible}; software is
    \textbf{logical and intangible}. The simple test is touch: if you can
    touch it, it is hardware.
  \item A \textbf{device driver} is \textbf{software}, not hardware. Its name
    follows the device, but it is a program (TRAP-10; PYQ-06).
  \item The \textbf{operating system} is \textbf{system software}, not
    application software (PYQ-06).
  \item \textbf{Firmware} is software stored in non-volatile memory inside
    hardware, such as the BIOS; it is software by nature, hardware by storage.
  \item \textbf{Humanware} (peopleware) is the \textbf{people} in the system,
    not a device or a program.
\end{enumerate}

\begin{example}
An item asks, ``Which of the following is software?'' with options CPU, a
printer, a device driver, and a keyboard. The answer is the \textbf{device
driver}: it is a program, whereas the CPU, the printer, and the keyboard are
physical. If the options instead are an operating system, MS Word, a BIOS
chip, and a game, the item is testing system software versus application
software and firmware, and the operating system is the system software.
\end{example}

\subsection*{Exercises}
\begin{enumerate}[label=\arabic*.]
  \item Define hardware and software, and give one example of each.
  \item Distinguish system software from application software. To which kind
    does the operating system belong?
  \item What is firmware, where is it stored, and why is it called ``between
    hardware and software''?
  \item Explain why a device driver is classified as software even though it
    is named after a hardware device.
  \item What is humanware? Give two roles it includes.
  \item List the layers of the computing stack from the bottom (hardware) to
    the top (application).
  \item Give one example each of internal hardware and external hardware.
  \item Using the print command as an example, explain how a request travels
    down the computing stack.
\end{enumerate}

\subsection*{Solutions}
\begingroup\small
\begin{enumerate}[label=\arabic*.]
  \item Hardware is the physical, tangible equipment of a computer, such as
    the CPU or a printer. Software is the programs and instructions that tell
    the hardware what to do, such as MS Word or a device driver.
  \item System software runs and manages the computer itself --- the operating
    system, device drivers, and utilities. Application software lets the user
    do a task, such as a word processor or a game. The operating system is
    \textbf{system software}.
  \item Firmware is software embedded in a hardware device and stored in
    non-volatile memory such as ROM or flash, for example the BIOS. It is
    software by nature but is fixed to a specific device, so it is described
    as sitting between hardware and software.
  \item A device driver is a program, so it is software. It is named after the
    device it serves, which tempts a candidate to call it hardware, but it is
    stored and works like other software; it merely translates the operating
    system's commands for one specific device.
  \item Humanware, or peopleware, is the human component of a computer system.
    It includes roles such as end users and system administrators (also
    programmers or operators).
  \item Hardware, then firmware, then device driver, then operating system,
    then application software.
  \item Internal hardware: the motherboard (also the CPU, RAM, or an internal
    drive). External hardware: the keyboard (also the mouse, monitor, printer,
    or an external drive).
  \item The user clicks Print in the application software. The application
    asks the operating system, which passes the job to the printer's device
    driver. The driver translates the command into the signals the printer
    hardware understands, and the printer's firmware controls the physical
    printing. The result returns up the stack to the user.
\end{enumerate}
\endgroup
\end{document}
